\subsection{Outer-Join Benchmark}
\label{sec:outerjoin-benchmark}

Standard analytical benchmarks do not isolate the factors that make outer-join
optimization difficult. TPC-H~\cite{tpch}, TPC-DS~\cite{tpcds}, JOB~\cite{job},
and SQLStorm~\cite{sqlstorm25} provide realistic schemas and join relationships,
but they do not systematically vary preservation semantics and join-tree shape.
We therefore use their queries as sources of realistic \emph{join skeletons},
rather than treating an inner-join query and its outer-join rewrite as
semantically equivalent workloads.

The core benchmark contains 34 traceable seeds: 5 from TPC-H, 7 from TPC-DS,
10 from JOB, and 12 from SQLStorm. From each seed, we extract a connected,
acyclic equality-join skeleton containing 3--6 relation instances. We retain
the relations and join predicates, normalize the surrounding query, and expose
attributes from multiple relations. Each generated selection is placed in the
global \texttt{WHERE} clause above the complete join tree, and every conjunct
is NULL-rejecting with respect to the relation it references. In particular,
we exclude NULL-tolerant guards such as $P(x)\lor x\ \texttt{IS NULL}$, as well
as \texttt{COALESCE} and \texttt{IS [NOT] DISTINCT FROM}. This restriction
keeps the performance workload within the null-rejection assumptions used by
OuterYan; NULL-tolerant predicates are retained only in separate correctness
tests.

The reported tree shapes and join vectors characterize the input SQL before
semantic simplification. A global NULL-rejecting predicate may legitimately
collapse an enclosing outer join. Such queries exercise the simplification
stage, while queries with surviving outer edges exercise outer-join enumeration
and execution.

\begin{figure}[t]
  \centering
  \small
  \begin{tabular}{c@{\;${}\longrightarrow{}$\;}c@{\;${}\longrightarrow{}$\;}c}
    \shortstack{34 source\\queries} &
    \shortstack{fixed connected\\join skeleton} &
    \shortstack{5 deterministic\\outer-join variants}
  \end{tabular}
  \caption{Benchmark construction. Each seed contributes the same five
  variants, producing a 170-query core.}
  \label{fig:outerjoin-benchmark-construction}
\end{figure}

Each skeleton produces five deterministic variants: a mixed left-deep tree
(\texttt{ldmx}); its semantics-preserving mirror, obtained by swapping operands
and LEFT/RIGHT join types (\texttt{rdmx}); a balanced bushy tree when the graph
permits one, otherwise a fixed alternative mixed tree (\texttt{bmx}/\texttt{alt});
and left- and right-deep uniform controls (\texttt{lduX} and \texttt{rduX}, where
$\mathsf{X}\in\{\mathsf{L},\mathsf{R},\mathsf{F}\}$). This gives 170 queries:
102 mixed and 68 uniform, with 80 left-deep, 68 right-deep, and 22 genuinely
bushy trees. Explicit parentheses fix the binary join tree because outer joins
are not freely associative or commutative.

This construction provides a controlled and representative measure of
outer-join performance. Holding the relations and join predicates fixed while
varying join type and tree shape isolates the system's handling of outer joins
from unrelated changes to query content. The combination of mixed and uniform
joins, left-deep, right-deep, and bushy trees, and multiple source schemas
covers the main semantic and structural factors that affect outer-join
optimization and execution. Projecting columns from multiple inputs makes
NULL-extension observable whenever it survives the declared NULL-rejecting
selection.

We also include the legacy workload in
\path{auto_test/outerjoin_query_old/} as a separate industry-oriented test
corpus. Its 170 outer-join queries are derived from workload patterns observed
in DBMS-X and instantiated over public benchmark schemas, thereby simulating
real industrial use cases while remaining reproducible. The corpus covers
practical patterns such as optional dimension enrichment, sparse relationships,
fact-completeness checks, and preserved master data, complementing the
controlled variants in the core benchmark.
